\begin{abstract}
Este artigo examina como relações intersetoriais conectam as emissões atribuídas pela produção e pela demanda no Brasil. Combina a MIP de 2015, com 67 atividades, a intensidades aproximadas de CO\textsubscript{2}, distinguindo volume, presença como fornecedor e concentração das atribuições. A rede é quase completa, mas seus pesos são desiguais. Energia distribui emissões entre diversos produtos finais; construção reúne emissões de outras origens e mantém elevada atribuição à própria atividade. O filtro de disparidade retém 13,40\% das ligações e 70,69\% do peso intersetorial, com cobertura desigual entre perfis. A análise identifica relações a investigar conjuntamente para mitigação, sem estimar efeitos de políticas. Os resultados ambientais são condicionais às intensidades utilizadas.
\end{abstract}
